
\newpage
~
\newpage

\includepdf[scale=1.020,pages=-]{others/Atm_Ocean.pdf}


\newpage

%%%%%%%%%%%%%%%%%%%%%%%%%%%%%%%%%%%%%%%%%%%%%%%%%%%%%%%%%%%%
\chapter{Manuscript description}


\poetry{13}{
Le climat, inconnue de demain,

\hspace*{1cm} est peut-être à saisir dans la méconnaissance que l'on a encore de l’océan.
}


\newpage


{\noindent\Large\textbf{The design process}}
\\


I wrote my thesis with the intention of making it accessible to individuals without a background in the field, enabling them to grasp some of its content and extract a meaningful message. Simultaneously, my aim was to provide knowledgeable readers with the opportunity to expand their understanding by discovering novel insights they might not have encountered before. Consequently, I viewed the process of writing this manuscript as a personal challenge. To achieve this, I realized that, after taking a several-month break following my thesis, I should be able to revisit all aspects of my research, including the complete framework and perspectives, simply by rereading my thesis. It was essential for me to avoid revisiting questions in the future that I had already explored and reflected upon. Therefore, I used this manuscript as a platform to organize and present my research as clearly and logically as possible. For example, in the introduction, I aimed to clarify what might have seemed poorly explained or difficult to understand during my thesis. In the conclusion and outlook, my intention was to share insights, even when certain aspects yielded no results, thus preventing us from pondering the same unanswered questions. I believed that this approach would make my thesis comprehensive and valuable for individuals who are curious, inquisitive, or new to the field, as well as for educational purposes. I hoped that by reading this manuscript, everything would make sense. That was my challenge, and only you, the readers, can tell me if I have succeeded. When I was writing my manuscript, I used to say: my delay matches my ambition, my progress matches my laziness.
\\

{\noindent\Large\textbf{My commitment}}
\\

I perceive commitment as a heartfelt response, one that arises from the depths, from a sense of responsibility. Because absorbing information triggers a bodily reaction in me, both in the moment, emotionally, and enduringly, passionately. Yet, I do not believe I am fundamentally different from others in this regard. However, as a climate scientist, facing climate change and the current research system, information on these subjects, I drink it in, I nourish myself with it, I breathe it; it is my daily life. So, it seeps in, spreads, and engulfs me. It was during this journey, throughout my thesis, that I encountered cognitive dissonance. Some find it charming and willingly embrace it. I, on the contrary, found it disheartening. It began to trouble my mind and gnaw at my body. I entered into an inner conflict, and a sense of responsibility began to incubate within me. So, it is high time I give birth to it, that I commit.

I have pondered extensively throughout this thesis about the nature of my commitment, and I believe that my position on this matter is not yet fully defined. I still have a reflective journey ahead to discover the path and the voice that suit me in this commitment process. Nevertheless, I saw this thesis as the perfect opportunity to solidify my commitment as it stands today, embracing its imperfections and inconsistencies, and giving it form within this manuscript by sharing a message. To structure it within a logical framework, one that has been shaped by my research during this thesis. I aspire to commit to raising awareness and igniting debates. To share, fundamentally, to exchange. Because, I believe, it is through this process that we make progress and construct a framework of coherent, just, and sustainable thoughts.

The ocean is at the heart of contemporary societal issues. Two Sustainable Development Goals (SDGs; \url{https://sdgs.un.org/goals}; last access: September 2023) are closely associated with it ('Climate action' and 'Life below water'). Currently, it is also the United Nations Decade of Ocean Science for Sustainable Development (2021-2030; \url{https://oceandecade.org/vision-mission/}; last access: September 2023) because the impact on the oceans poses a growing threat to marine ecosystems. The vision proposed for this decade is: "The science we need for the ocean we want". What a beautiful statement! It encapsulates everything: science, the ocean, an opportunity to reflect, question, and explore what lies beyond these words.

As you will discover throughout this manuscript, some scientists were already advocating for a sensitive approach to the climate issue in the past, with many of them contemplating the path to follow. After all, a scientist remains sensitive, fundamentally human, experiencing emotions, and wrestling with passions, which often, as in my case, lead to be passionate. Sensitivity is, in fact, intertwined with a set of values that are uniquely our own:

\begin{spacing}{1.0}
~
\end{spacing}

\centerline{\begin{minipage}{0.85\textwidth}
\setlength\parindent{12pt}
\textit{It is an illusion to think that a researcher can completely divorce themselves from their values: every scientific pursuit is a human endeavor, embedded in a social context, and thus influenced by values. Therefore, the primary challenge is not to expect researchers to be devoid of values but to have those values explicitly stated and to adhere to the principles of integrity and rigor that should characterize scientific inquiry.\footnote{Il est illusoire de penser que le chercheur puisse se débarrasser totalement de ses valeurs : toute science est une entreprise humaine, inscrite dans un contexte social et, ce faisant, nourrie de valeurs. L’enjeu premier n’est donc pas d’attendre du chercheur qu’il en soit dépourvu mais qu’il les explicite et qu’il respecte les exigences d’intégrité et de rigueur qui doivent caractériser la démarche scientifique.}} \citep{comets__2023}
\end{minipage}}

\begin{spacing}{1.0}
~
\end{spacing}

In fact, I contemplate whether a researcher's role should also encompass breathing life into their research. Why delegate the responsibility of infusing vitality into the outcomes of our research to others when no one understands the subject better than the person who generated the results? Breathing life into one's research is synonymous with infusing it with and sensitivity. My primary objective is to breathe life into it through words, as words set in motion and establish the tone. Consequently, it becomes the responsibility of each individual to be discerning in listening to this message, capable of experiencing emotions, and subsequently distinguishing the objectified content from the subjectified emotions.



Here, I am challenging our research practices and the dialogue with citizens. However, this commitment does not represent a normative position, which is the focus of the recent opinion from the CNRS Ethics Committee (COMETS) titled 'Between Freedom and Responsibility: The Public Engagement of Researchers'\footnotemark{}\footnotetext{'\textit{Entre liberté et responsabilité : l'engagement public des chercheurs et des chercheuses}'.}. It is essentially an ethical approach that drives me to engage through this manuscript, a perspective further contextualized by the words of \citet{coutellec_penser_2019}:

\begin{spacing}{1.0}
~
\end{spacing}

\centerline{\begin{minipage}{0.85\textwidth}
\setlength\parindent{12pt}
\textit{
At the very least, research ethics can be described as a reflective approach to the values and purposes of scientific research; scientific integrity as a normative approach aimed at regulating (good) practices within a community by establishing norms and principles; and the social responsibility of sciences as a political approach that seeks to understand the context and anticipate the consequences of science, with an awareness of its involved nature.
\newline \indent
[...]
\newline \indent
The ideal approach to ethical questioning by young researchers should extend beyond merely alternating between case studies on their practices and the vertical transmission of a catalog of recommendations from codes, charters, or guides of good practices. The real challenge is to instill a culture of ethical inquiry, allowing the integration of this dimension at the core of research work, irrespective of the discipline.
\newline \indent
[...]
\newline \indent
Such training could thus serve as opportunities to take a step back on research and sciences, reflecting on the nature of the knowledge produced and its societal implications.
\newline \indent
[...]
\newline \indent
to position research ethics as one of the cornerstones in transforming our relationship with science and, more broadly, our connection with knowledge [...].\footnote{De façon minimale, il est possible de qualifier l’éthique de la recherche comme démarche réflexive sur les valeurs et les finalités de la recherche scientifique ; l’intégrité scientifique comme démarche normative qui vise à encadrer les (bonnes) pratiques d’une communauté, en établissant normes et principes ; la responsabilité sociale des sciences comme démarche politique qui vise à appréhender le contexte et anticiper les conséquences de la science dans une prise de conscience du caractère impliqué de celle-ci.
\newline \indent
[...]
\newline \indent
La démarche souhaitable d’appropriation du questionnement éthique par les jeunes chercheurs ne peut se résumer à une alternance entre étude de cas sur leur pratique et transmission verticale d’un catalogue de recommandations issues de codes, de chartes ou de guides de bonnes pratiques. L’enjeu est plutôt l’acquisition d’une culture du questionnement éthique permettant d’intégrer cette dimension au coeur du travail de recherche, quelle que soit la discipline.
\newline \indent
[...]
\newline \indent
De telles formations pourraient ainsi être des occasions d’une prise de recul sur la recherche et les sciences, sur la nature des connaissances produites et sur la portée de celle-ci en société.
\newline \indent
[...]
\newline \indent
faire de l’éthique de la recherche l’un des piliers d’une transformation de notre rapport à la science et, plus généralement, de notre rapport au savoir [...].}
}
\end{minipage}}

\begin{spacing}{1.0}
~
\end{spacing}

These opening lines of my commitment, to me, serve as a way to contextualize it, to define the ethical framework that encompasses both my freedom of expression and my social responsibility, without in any way compromising the scientific integrity of this manuscript. In essence, moving beyond a purely scientific pursuit, conducted with a scientific approach as objective as possible, I infuse life into the message it encapsulates through a sensitive and reflective approach. To maintain transparency, I would like to mention the sections in which I express my commitment, initially focused on climate sciences and the broader Anthropocene, and subsequently extended to the research system: (i) in Chapt.~\ref{chap1}, the 'Preamble', Box I.1, and 'Opening reflections'; (ii) in Chapt.~\ref{chap6}, the 'Preamble', Sect. VI.3 and 'Closing reflections'; and (iii) in Sect.~\ref{annex_temoignage} as well as Sect.~\ref{annex_nouvelle} (in French though). To me, commitment equates to assuming one's responsibilities. This commitment, which I share, carries the weight of my sole responsibility.
\\





{\noindent\Large\textbf{Structural elements}}
\\

I also aimed to make my thesis visually appealing while conveying a message. As you may have noticed on the front and back covers, you will find the famous 'climate stripes' (\url{https://showyourstripes.info}; last accessed in September 2023), providing a visual representation of the global surface temperature anomaly relative to the 1971-2000 period. Blue signifies a negative anomaly, while red indicates a positive one. These climate stripes span the period from 1850 to 2022, and they are read from bottom to top.

Furthermore, at the bottom of each page, you will find stripes representing the global evolution of the atmospheric CO\textsubscript{2} concentration (at the bottom of the left-hand pages) and the surface ocean \textit{p}H (at the bottom of the right-hand pages). These data are linearly interpolated from annual values. For CO\textsubscript{2} concentrations, we used values prescribed in the CMIP6 exercise for the various scenarios considered, and for the surface ocean \textit{p}H, we derived them from the IPSL model (IPSL-CM6A-LR) for the respective scenarios (see Sect.~\ref{sect:These_Method_CMIPs}). The historical period, gradually displayed in the first part of the thesis, covers the years from 1850 to 2014. In contrast, the scenarios, progressively shown in the second part, encompass the period from 2015 to 2100. Please note that I adjusted the scale after the end of the historical period so that we can discern trends both during the historical period and throughout the scenarios. I performed temporal interpolation to ensure coverage of the entire historical period and scenarios throughout the thesis. For the period spanning from 2015 to 2100, for improved readability, I have chosen to display only two extreme scenarios: SSP1-2.6 and SSP5-8.5, corresponding to low-emissions and high-emissions scenarios (see Box I.3), which are presented in parallel at the top and bottom, respectively. Lastly, during the historical period, we included actual measurements of the atmospheric CO\textsubscript{2} concentrations obtained from the Mauna Loa observatory in Hawaii (USA) as soon as they became available.

This thesis is a reflection of my personality, incorporating a sensitive approach expressed through poetic interludes placed between each section. These interludes are presented in French only, as I do not presume to be able to translate them with the subtleties they contain. They resonate with my passion for writing, offering a sensory exploration of words from my perspective. Specifically, in Sect.~\ref{annex_nouvelle}, I share a short story written during this thesis as part of a creative climate writing workshop.

Finally, this thesis includes watercolor illustrations that complement specific messages. These artworks were skillfully crafted by my father, and I want to express my heartfelt gratitude and admiration for his invaluable contributions. Thank you, Dad!